\clearpage
\fancyhead[L]{\small PHILOSOPHICAL TRANSACTIONS.}
\newcommand{\ObservedWollastonDenominator}{\hbox{\vrule width0pt height4.600000bp depth0pt\special{pdf:bcontent}\special{pdf:code q 0 g 2.453333 4.600000 m 2.146667 4.293333 l 2.453333 3.986667 l 2.453333 3.373333 l 1.840000 2.760000 l 1.226667 2.760000 l 1.226667 2.453333 l 0.613333 1.840000 l 0.306667 1.840000 l 0.000000 1.533333 l 0.000000 0.613333 l 0.613333 0.000000 l 3.373333 0.000000 l 3.986667 0.613333 l 3.986667 0.920000 l 4.293333 1.226667 l 4.293333 1.533333 l 3.986667 1.840000 l 3.986667 2.146667 l 3.680000 2.146667 l 3.373333 2.453333 l 3.373333 3.066667 l 3.680000 3.373333 l 3.680000 3.986667 l 3.066667 4.600000 l h 1.226667 2.146667 m 1.533333 2.453333 l 1.840000 2.453333 l 2.453333 1.840000 l 2.453333 1.226667 l 1.840000 0.613333 l 1.226667 1.226667 l h f* Q}\special{pdf:econtent}\kern4.293333bp}}

\hypertarget{native-566-lower}{}
\begin{center}\small Göttingische gelehrte Anzeigen. No. 27. Pages 269..272. 14 February 1814.\\[3pt]\rule{.72\linewidth}{.3pt}\end{center}

\emph{Philosophical Transactions of the Royal Society of London} for the Jear 1813. IV and 304 pages. Meteorological Journal and 8 pages of index, in quarto.

\emph{Observations of the summer solstices of} 1812 \emph{and} 1813, \emph{and of the winter solstice of} 1812, \emph{at the Royal Observatory at Greenwich}, by J. \textsc{Pond} (three papers). The instrument used for these observations is a six-foot mural circle made by \textsc{Troughton}, of which Mr. \textsc{Pond} promises to give a detailed description later. The few details given here show that its distinctive feature is that it measures arcs of the celestial meridian without reference to the zenith. It therefore has neither a plumb line nor a level. The determination of the position of the zenith in the meridian is left solely to the zenith sector. The accuracy claimed here for measurements with this instrument, which is said to cost 1600 guineas according to a private report received via Paris, seems to promise a new era for this part of practical astronomy. The table gives the results of these determinations for the mean obliquity of the ecliptic.

\begin{center}\fontsize{9}{11}\selectfont\setlength{\tabcolsep}{3pt}\renewcommand{\arraystretch}{1.18}\begin{tabular}{@{}l c r@{\;}r@{\;}r@{}}
Summer solstice of 1812 & $.\;.\;.\;.\;.\;.\;.\;.\;.\;.$ & $23^\circ$ & $27'$ & $50''5$\\
Winter solstice of 1812 & $.\;.\;.\;.\;.\;.\;.\;.\;.\;.$ & $23$ & $27$ & $47.35$\\
Summer solstice of 1813 & $.\;.\;.\;.\;.\;.\;.\;.\;.\;.$ & $23$ & $27$ & $49.5$\\
\end{tabular}\end{center}

{\rightskip=0pt plus 3em In calculating the observations, Mr. \textsc{Pond} used \textsc{Bradley}'s refraction tables.\par}

\emph{Catalogue of the north polar distances of} 44 \emph{of the principal fixed stars}, for the beginning of 1813, and in a later paper \emph{a similar catalogue for} 84 \emph{fixed stars}, by J. \textsc{Pond}. The first of these \SourceLeaf{567}catalogues had already become known in Germany too. The second is to be regarded as an enlarged and improved edition. The number of observations on which these determinations are based reaches 100 for some stars and even 200 for Polaris. For some of the principal stars Mr. \textsc{Pond} also gives the results of the individual series of observations (ten observations in each). Their agreement attests to the excellence of the instrument and surpasses anything of the kind known so far. For many stars Mr. \textsc{Pond} considered his determination reliable to $\frac{1}{4}$ of a second, so far as it depends solely on the instrument. Mr. \textsc{Pond} has also compiled a selection of 30 of the principal stars north of the equator in a separate list that he calls the standard catalogue (Standard catalogue). He intends to keep improving it as a counterpart to \textsc{Maskelyne}'s catalogue for right ascensions.

Two papers by \textsc{William Hyde Wollaston} also relate to practical astronomy. The first gives a method of drawing extremely fine metal wires. It consists in placing a gold or platinum wire inside an initially thick silver wire (either by boring out the silver wire or by casting molten silver around the platinum wire in a mould), drawing out this double wire in the usual way, and finally dissolving away the silver with warm nitric acid. These experiments succeeded best with platinum. In this way he obtained threads $\frac{1}{1000}$ and $\frac{1}{8000}$ of an inch thick that were very well suited to being stretched at the focus of astronomical telescopes. These threads are already far thinner than the finest spider threads. (The reviewer found by measurement that the thickness of a spider thread stretched at the focus of a 16-inch telescope covered an arc of $7''8$. The thickness of that thread would accordingly be $\frac{1}{1630}$ of an inch.) Mr. \textsc{Wollaston}, however, went much further and produced threads of almost incredible fineness (though only in short pieces). They were still quite strong in proportion to their size. A platinum thread $\frac{1}{18000}$ of an inch thick could still bear $1\frac{1}{\ObservedWollastonDenominator}$ grains. To measure the thickness of such exceedingly delicate objects with a certain degree of accuracy, he used a special device that he describes in the other paper under the name Single-lens micrometer. Its essential part is a lens of very short focal length ($\frac{1}{12}$ of an inch), with the object placed at its focus. The aperture of this lens is so small that a small hole can be made beside the lens in the brass plate in which it is mounted, at a distance of $\frac{1}{23}$ of an inch from the centre. The pupil receives light simultaneously from the object through the lens and from a movable scale through the hole. The scale is moved away until the object exactly covers one or a few of its divisions. The distance of the scale and the number of divisions then determine the size of the object.

\Needspace{26\baselineskip}
Also relevant here is \emph{On the light of the Cassegrain telescope compared with that of the Gregorian}, by H. \textsc{Kater}. Two pairs of such telescopes, made by a single instrument maker (Mr. \textsc{Crickmore} in Ipswich), were compared. The comparison took into account the difference in the areas of the large mirrors and the loss caused by the smaller mirrors and the arms that supported them. In one pair the ratio of the amounts of light was 7 to 3, in the other 3 to 2, both times in favour of the \textsc{Cassegrain} telescope. In the first pair the mirror of the \textsc{Cassegrain} had the advantage of a somewhat better polish, and in the other pair the mirror of the \textsc{Gregory} had that advantage. The author therefore believes that under exactly equal conditions the \textsc{Cassegrain} telescope has about twice as much light as the \textsc{Gregory}. If we may assume that the opening at which the eye was placed directly was of the proper size in all four telescopes to let all the light from the large mirrors pass through — the author does not touch on this point, but the reviewer \SourceLeaf{568}has handled several \textsc{Gregory} telescopes in which this opening was too small — then these findings are of the greatest interest to both the astronomer and the physicist. One might be inclined to accept the author's explanation. He seeks the reason for the curious phenomenon in the fact that two real images are formed in the \textsc{Gregory} telescope, whereas only one is formed in the \textsc{Cassegrain}, and infers from this that the crossing of the rays at a point, or the formation of a real image, causes a certain loss of light. He thinks this may be connected with the experience that \textsc{Galilei}’s telescopes are more effective than astronomical telescopes under the same conditions.
